%! Conic Sections — Unit 8, Lesson 8.5 — Exit Ticket (Answer Key)
\documentclass[10pt]{article}
\usepackage{precalculus-article}
\usepackage{precalculus-key}   % pulls in -boxes; do NOT also load -boxes
\newcommand{\TallMath}[1]{$\displaystyle #1\rule[-1.4em]{0pt}{3.2em}$}

\begin{document}

\pageheader{Unit 8, Lesson 8.5}{Exit Ticket --- Answer Key}
\namedateperiod

\vspace{0.12in}

\begin{notesbox}{Conic Sections Check}

\textbf{1.}\quad Write the standard form equation of an ellipse centered at
$(3, -1)$ with horizontal radius $5$ and vertical radius $2$.

\ansline{$\dfrac{(x-3)^2}{25} + \dfrac{(y+1)^2}{4} = 1$}

\vspace{0.06in}

\textbf{2.}\quad For the equation $\dfrac{(y+2)^2}{16} - \dfrac{(x-1)^2}{9} = 1$:

\vspace{4pt}
(a)\quad What type of conic is this?
\ansline{Hyperbola (the $y$-term is positive and the two terms are subtracted)}

(b)\quad What is the center?
\ansline{$(1,\,-2)$}

(c)\quad Write both asymptote equations.

\ansline{$y + 2 = \dfrac{4}{3}(x-1)$ \quad and \quad $y + 2 = -\dfrac{4}{3}(x-1)$}

\end{notesbox}

\begin{teachernote}
Problem 1: accept any algebraically equivalent form. The most common error is
writing $25$ and $4$ without squaring (i.e., $a = 5 \Rightarrow a^2 = 25$) ---
verify students wrote $a^2$ in the denominator, not $a$.\\[3pt]
Problem 2: $b^2 = 16 \Rightarrow b = 4$ (from the positive $y$ term) and
$a^2 = 9 \Rightarrow a = 3$ (from the $x$ term). Asymptote slope is $b/a = 4/3$.
The hyperbola opens up and down because the $y$-term is positive.
\end{teachernote}

\end{document}
